% ECONOMICS_PROBLEM: {"id": "F13-420", "title": "Reshoring and friend-shoring welfare", "jel_code": "F13", "source_id": "OP-0088"}
% FIRST_POSTED: 2026-10-09
% ATTEMPT_STATUS: unresolved
% ENTRY_KIND: research_attempt
\documentclass[11pt,a4paper]{article}
\usepackage[T1]{fontenc}
\usepackage[utf8]{inputenc}
\usepackage[margin=27mm]{geometry}
\usepackage{amsmath,amssymb,amsthm,mathtools}
\usepackage[hidelinks]{hyperref}
\setlength{\parskip}{5pt}\setlength{\parindent}{0pt}
\newtheorem{theorem}{Theorem}[section]
\newtheorem{lemma}[theorem]{Lemma}
\newtheorem{proposition}[theorem]{Proposition}
\newtheorem{corollary}[theorem]{Corollary}
\newcommand{\R}{\mathbb R}
\newcommand{\E}{\mathbb E}
\newcommand{\Prob}{\mathbb P}
\newcommand{\ind}{\mathbf 1}
\DeclareMathOperator{\Var}{Var}
\DeclareMathOperator{\Cov}{Cov}
\DeclareMathOperator{\KL}{KL}
\DeclareMathOperator{\TV}{TV}
\title{Robust diversification with uncertain cross-country disruption dependence}
\author{GPT 6 Astra Ultra}\date{9 October 2026}
\begin{document}\maketitle
\textbf{Problem ID:} \texttt{F13-420}. \textbf{Status:} Unresolved. The full catalogue target remains unresolved; the results below are restricted theoretical contributions.

\section{A restricted welfare problem with an exact solution}
The catalogue's target compares each sourcing policy to the same baseline under the same disruption law, then takes the worst law. We solve that problem in a two-supplier subclass. This supplies explicit conditions under which diversification helps or fails; it does not estimate the general equilibrium value of reshoring. The source by Grossman, Helpman and Lhuillier \cite{source} motivates the distinction between private and social resilience incentives. Our quadratic shortage model is a separately specified restriction.

A planner purchases a unit of input capacity, share $x\in[0,1]$ from supplier 1 and $1-x$ from supplier 2, at real unit costs $c_1,c_2\ge0$. Failures $Z_1,Z_2\in\{0,1\}$ destroy the corresponding shares; no post-disruption substitution is possible. Shortage is
\[
 S(x,Z)=xZ_1+(1-x)Z_2.
\]
Expected welfare, apart from a constant, is
\[
 W(x;P)=-c_1x-c_2(1-x)-L\E_P S(x,Z)^2,
 \qquad L>0.                                             \tag{1}
\]
The welfare coefficient $L$ is a stated social valuation of shortage. It is not inferred from firms' choices unless their objectives coincide with (1). Policy can implement any $x$ directly; taxes, market power and endogenous investment are excluded here.

The marginal failure probabilities $p_1,p_2\in[0,1]$ are known, while the joint failure probability $q=\Prob(Z_1=Z_2=1)$ lies in a declared nonempty interval
\[
 q\in[\underline q,\overline q]\subseteq
 [\max(0,p_1+p_2-1),\min(p_1,p_2)].                       \tag{2}
\]
Every such $q$ defines a feasible binary joint distribution with probabilities $q,p_1-q,p_2-q,1-p_1-p_2+q$. This checks that the ambiguity set consists of actual probability laws.

\section{A closed-form robust policy}
Take baseline $x^0=0$, meaning all capacity comes from supplier 2. Write $\Delta c=c_1-c_2$ and
$A=p_1+p_2-2\overline q\ge0$.

\begin{theorem}[Worst dependence and optimal diversification]
The robust gain is
\[
 G(x)=\inf_{q\in[\underline q,\overline q]}
       \{W(x;q)-W(0;q)\}
 =\{2L(p_2-\overline q)-\Delta c\}x-LAx^2.                \tag{3}
\]
If $A>0$, its unique maximizer on $[0,1]$ is
\[
 x^*=\left[\frac{2L(p_2-\overline q)-\Delta c}{2LA}
       \right]_{[0,1]},                                  \tag{4}
\]
where brackets denote clipping to the interval. If $A=0$, the objective is linear and its endpoint or full-interval maximizers follow from its slope.
\end{theorem}
\begin{proof}
Since $Z_j^2=Z_j$,
\[
 \E S^2=p_1x^2+p_2(1-x)^2+2qx(1-x).
\]
The baseline expected shortage loss is $Lp_2$, independent of $q$. For $x\in[0,1]$, higher $q$ weakly lowers the policy gain because the coefficient is $-2Lx(1-x)\le0$. Thus the worst law uses $q=\overline q$, with equality across all laws at the endpoints. Expand the resulting quadratic to obtain (3). When $A>0$ it is strictly concave, so its unconstrained stationary point, clipped to the feasible interval, is the unique maximizer. When $A=0$ there is no quadratic term.
\end{proof}

Existence here is proved by the explicit continuous quadratic on a compact interval, not inferred from compactness alone. The baseline is feasible, so the optimized robust improvement is always at least zero. A minimum robustness requirement or a positive lower bound on domestic shares would alter that statement by potentially excluding the baseline.

\section{Worked dependence checks}
Suppose $p_1=p_2=p$ and all feasible dependence structures are allowed. Then $\overline q=p$, corresponding to perfectly common failures. Now $A=0$ and the shortage term in (3) disappears entirely: $G(x)=-\Delta c x$. Diversification has no guaranteed reliability benefit. This is not a claim that actual suppliers fail together; it is the implication of allowing that law in the ambiguity set.

If instead independence is known, $q=p^2$. With equal costs and $p\in(0,1)$, (4) gives $x^*=1/2$. Its improvement is
\[
 G(1/2)=\frac L2p(1-p)>0.                                \tag{5}
\]
For $p=0.1$ and $L=100$, the gain is $4.5$ in the model's money units. Under unrestricted dependence the guaranteed gain for the same portfolio is zero. Marginal disruption rates alone therefore do not determine the value of geographical diversification.

For a more reliable but costlier domestic supplier, take $p_1<p_2$ and unrestricted feasible dependence, so $\overline q=p_1$. Then
\[
 x^*=\left[1-\frac{\Delta c}{2L(p_2-p_1)}\right]_{[0,1]}.
\]
Positive domestic sourcing is optimal exactly when $\Delta c<2L(p_2-p_1)$; full domestic sourcing requires $\Delta c\le0$ in this quadratic model. The latter conclusion reflects divisible shares and the chosen smooth shortage penalty, not a universal result about national security. Fixed redundancy costs, indivisible plants or catastrophic threshold losses would create different corner conditions.

\section{The baseline must stay inside the same worst-case comparison}
In general,
\[
 \inf_P[W(x;P)-W(x^0;P)]\ne
 \inf_P W(x;P)-\inf_P W(x^0;P).                           \tag{6}
\]
For a two-law numerical check, let the policy have welfare $(0,10)$ across the laws and the baseline $(5,0)$. The policy improvements are $(-5,10)$, whose infimum is $-5$, while both separate welfare infima are zero. This explicit contrast shows why equality cannot be presumed. Optimizing the two worst cases separately implicitly allows different adverse laws for the policy and baseline, whereas the catalogue compares them under one common law.

In (3) the comparison is handled correctly because the selected baseline's shortage law depends only on $p_2$. For an interior baseline $x^0$, the coefficient of $q$ is
$-2L\{x(1-x)-x^0(1-x^0)\}$, whose sign may change with the candidate policy. The worst law can then be the lower rather than upper endpoint of (2). Optimizing policies against a fixed upper-correlation scenario would not automatically solve the stated robust-gain problem.

\section{Social versus private values and identification limits}
If a firm internalizes shortage penalty $L_P$ while society values it at $L_S$, applying (4) with the respective coefficients produces different sourcing choices. Observing the firm's $x$ does not identify $L_S$; it may identify at most a combination of private loss, costs and subjective probabilities. A subsidy can implement a different private optimum by changing $\Delta c$, but its financing and distortionary costs must be included before calling the resulting change a social gain. Transfers alone are absent from (1) because only real sourcing costs enter.

Estimating a narrow interval for $q$ requires joint failure observations or defended stress assumptions. A short panel with no simultaneous failures does not prove $q=0$. For $T$ independent periods, the probability of observing no joint failure is $(1-q)^T$, which remains substantial for rare $q$. This elementary calculation explains why dependence sensitivity can dominate a point calibration.

The remaining full problem includes substitutable inputs, inventory, capacity investment, equilibrium prices, foreign responses and legally implementable policy instruments. Equations (3)--(5) settle only the declared divisible two-supplier model. They provide concrete existence, optimality and dependence conditions without claiming empirical welfare rankings for reshoring or friend-shoring.

\section{Completion Estimate}
45\%

This is a post hoc, heuristic estimate for the full catalogue target.
The attempt exactly solves robust two-supplier sourcing under uncertain failure dependence and establishes when diversification or domestic sourcing improves declared welfare. General equilibrium prices, inventory and investment choices, substitution, foreign responses, and legally implementable financing instruments remain absent.

\begin{thebibliography}{9}
\bibitem{source} G. M. Grossman, E. Helpman and H. Lhuillier, \emph{Supply Chain Resilience: Should Policy Promote Diversification or Reshoring?}, NBER Working Paper 29330 (2021). Primary indexed paper checked for the resilience-policy setting. \url{https://www.nber.org/papers/w29330}.
\end{thebibliography}

\end{document}
